\chapter{Software System \& Experiment Infrastructure}
\label{ch:software}

The reproducibility claims made in Chapter~\ref{ch:sota-repair} are only meaningful if they
are checkable, and the finding that a published baseline's error bars were misattributed
(Section~\ref{sec:btgt-reproducibility}) was itself only discoverable because every run left
a complete, self-describing artifact behind. This chapter documents the infrastructure that
made that possible. It is presented as a contribution rather than as an implementation
appendix, because the central methodological result of this thesis could not have been
obtained without it.

%==============================================================================
\section{Registry-Driven Experiment Execution}
\label{sec:runner}

Experiments are not invoked as ad-hoc scripts. Each is a declarative entry in a YAML
registry --- identifier, mode, model, dataset, seed, notebook, adapter, configuration path,
checkpoint path, threshold mode --- and a single runner executes them via papermill.

\begin{quote}\small\ttfamily
python run\_experiment.py --id recon-ablation-gelstm-none \\
python run\_experiment.py --id recon-ablation-gelstm-none --dry-run \\
python run\_experiment.py --status \quad|\quad --collect
\end{quote}

The consequences that matter for this thesis:

\begin{description}[leftmargin=2.8cm,style=nextline]
  \item[One place per knob] An ablation arm differs from its reference by exactly the field
    that names the arm. Because the shared hyperparameters live in one configuration file
    referenced by every arm, ``everything else was identical'' is enforced by construction
    rather than asserted.
  \item[\texttt{resolved\_config.json}] Every run writes the fully merged parameter set it
    actually used. Verifying that four ablation arms differ only in \texttt{encoder\_init}
    is a diff, not an act of faith.
  \item[Git provenance] Every run records the commit it ran at and whether the tree was
    dirty. Pooling runs across a code change is therefore detectable rather than silent.
  \item[\texttt{--dry-run}] Prints the merged parameters and exits, so a misconfigured sweep
    is caught before it consumes GPU time rather than after.
  \item[A results ledger] \texttt{--collect} rebuilds a single \texttt{RESULTS.csv} across
    every run of every experiment. Appendix~C is generated from it rather than transcribed.
\end{description}

\edge{The reproducibility finding was a consequence of the infrastructure, not of
suspicion.}{%
Seven completed runs of one Brain-TokenGT configuration at one seed sat in the output tree
because the runner never overwrites a run directory --- each execution gets its own,
timestamped and named. A workflow that overwrites, or that keeps only the latest, would have
destroyed exactly the evidence that revealed the non-determinism. The design choice that
made this visible was made for other reasons, which is worth stating: reproducibility
infrastructure pays out in ways that cannot be anticipated when it is built.}

%==============================================================================
\section{Layered Architecture}
\label{sec:layering}

A strict separation is enforced between four layers:

\begin{itemize}
  \item \texttt{model/} --- pure modelling logic. No file I/O, no experiment-tracker calls,
    no path construction, and no hardcoded hyperparameters.
  \item \texttt{configs/} --- typed dataclasses with explicit defaults. A new
    hyperparameter is a new field, never a loose keyword argument threaded through call
    sites.
  \item \texttt{notebooks/} --- orchestration only: load data, build a config, call into
    \texttt{model/}, persist artifacts. Training loops do not live in notebooks.
  \item \texttt{common/} --- shared utilities with no model-specific imports: seeding,
    splits, cross-validation, sanity audits, calibration, confound diagnostics.
\end{itemize}

This is what made the four-arm ablation a one-field change rather than a fork of the
training code, and what allows the \texttt{none} arm --- which builds no encoder and widens
the recurrent input from 64 to 200 --- to be expressed without a parallel code path.

%==============================================================================
\section{Failing Loudly by Design}
\label{sec:fail-loud}

A recurring principle in this codebase is that a silent fallback is worse than a crash. Two
instances are load-bearing for the results:

\begin{enumerate}
  \item The classifier evaluation entry point \textbf{raises} when no explicit decision
    threshold is supplied, rather than defaulting to $0.5$. The error message names the leak
    it is preventing (Section~\ref{sec:thresholds}).
  \item Setting \texttt{encoder\_init} and the legacy \texttt{freeze\_encoder} flag to
    contradictory values raises rather than silently resolving one of them --- so a
    half-migrated configuration cannot quietly run the wrong ablation arm.
\end{enumerate}

The counter-example that motivates the principle is documented in
Section~\ref{sec:external-cohorts}: a visit parser that returned \texttt{None} instead of
raising produced an empty ADNI cohort with no error at all.

%==============================================================================
\section{Testing and Continuous Integration}
\label{sec:ci}

A test suite runs on every push and pull request, alongside linting. Type checking, format
checking, complexity and security analysis run as \emph{ratcheted} non-blocking checks: a
pre-existing backlog is tolerated, but a change that introduces a new finding fails the
build. This keeps quality monotonically improving without demanding a repository-wide
cleanup before any new work can land.

The tests that matter most here are the ones pinning \emph{no-op-by-default} contracts ---
that a newly added flag leaves existing behaviour byte-identical when unset. The
\texttt{encoder\_grad} flag is covered this way, because without it the \texttt{random}
ablation arm would silently have measured a frozen encoder
(Section~\ref{sec:encoder-arms}).

%==============================================================================
\section{Dual-Host GPU Dispatch}
\label{sec:dispatch}

Training runs across two workstations that mount the same network export at the same path,
so the repository, the virtual environment, the data and the outputs form a single shared
tree. A dispatch utility assigns experiments to whichever GPU has more free memory and can
run a sweep across both boxes simultaneously.

The shared tree gives free parallelism and no mutual exclusion, which is a real hazard
rather than a theoretical one: the runner rewrites each experiment's \texttt{latest} pointer
without locking, so the same experiment identifier started twice --- on one box or across
both --- leaves that pointer aimed at whichever run finished last, and downstream notebooks
silently read the wrong run. Work is therefore split across hosts \emph{by experiment
identifier}, never by running one identifier in two places, and the runner enforces this
before creating any run directory. Liveness of a run owned by the other host is established
from a heartbeat file on the shared tree rather than from a process table, and a run whose
liveness cannot be established is reported as unknown and treated as alive --- because
marking a healthy run dead corrupts it for every host.

%==============================================================================
\section{Interactive Research Dashboard}
\label{sec:dashboard}

A FastAPI backend with a Vite front end supports cohort browsing, single-subject trajectory
inspection, and inference against a trained trajectory model. It reuses the project
environment for the model path rather than duplicating the deep-learning dependency stack.

\draftnote{Scope check before this section is finalised: describe only what is deployed and
  demonstrable. If the inference path is not wired end-to-end to a current checkpoint, say
  so here and present the dashboard as cohort exploration rather than as a clinical decision
  support prototype.}
